\documentclass[a4paper,12pt]{article}
\usepackage[utf8]{inputenc}
\usepackage{graphicx}
\usepackage{float}
\usepackage{amsmath}
\usepackage{caption}
\usepackage{geometry}
\geometry{margin=1in}

\title{Digital Logic Design Report}
\author{Mahla Entezari}
\date{Academic Year 2023-2024}

\begin{document}

\maketitle

\section*{Instructor: Shima Tabibian}
\section*{Course: Principles of Computer Systems Laboratory}

\section{Binary to Gray and Gray to Binary Code Converter}
In this project, a converter was implemented to transform a 4-bit binary number into its corresponding Gray code and vice versa. The conversion formulas used are as follows:

\begin{itemize}
    \item \textbf{Binary to Gray:} The most significant bit (MSB) remains unchanged. Each subsequent bit is obtained by XORing the current binary bit with the previous one.
    \item \textbf{Gray to Binary:} The MSB remains unchanged. Each subsequent binary bit is found by XORing the previous binary bit with the current Gray code bit.
\end{itemize}

\begin{figure}[H]
    \centering
    \includegraphics[width=0.9\textwidth]{1.png}
    \caption{Binary to Gray and Gray to Binary Conversion Circuit}
\end{figure}

The circuit was tested successfully with two different 4-bit values to verify accurate conversion in both directions.

\section{Full Adder Design}
The full adder circuit was designed and implemented using basic logic gates. The full adder adds two input bits and a carry-in bit, producing a sum and a carry-out bit.

\begin{figure}[H]
    \centering
    \includegraphics[width=0.9\textwidth]{2.png}
    \caption{Full Adder Circuit Designs}
\end{figure}

The circuit was tested across all 8 possible combinations of 3 input bits (A, B, and Carry-in), verifying the correct sum and carry-out results.

\section{3-bit Binary Adder Using Full Adders}
A 3-bit binary adder was constructed by chaining three full adders in series. Each full adder handles a corresponding bit from the input binary numbers.

\begin{figure}[H]
    \centering
    \includegraphics[width=0.9\textwidth]{3.png}
    \caption{3-bit Full Adder Circuit}
\end{figure}

This design was tested with different 3-bit binary numbers to validate the overall addition result and the final carry-out.

\section{Conclusion}
In this laboratory project, we explored various fundamental digital logic circuits:

\begin{itemize}
    \item Code converters (Binary \textless$\rightarrow$\textgreater Gray)
    \item Full adder circuits
    \item Multi-bit binary adders using chained full adders
\end{itemize}

Each design was implemented and simulated successfully using Logicly software. These circuits form the foundation of digital arithmetic units in computer systems.

\end{document}

